%%==========================================================================
%% This is the very first version of a Streetwise LaTeX template. By
%% chosing the document class, it is easily convertible to a two-column IEEE
%% paper format, provided any conflicting packages and (new)commmands are
%% removed.
%%
%% Erwin de Gelder, January 6, 2018
%%==========================================================================

%%==========================================================================
%% Title:
%% Author(s):
%% To be published in:
%% Date of submission:
%% Date of submission R1:
%% Date of submission R2:
%% Date of final submission:
%%==========================================================================


%% Generic
\documentclass[10pt,final,a4paper,oneside,onecolumn]{article}

%% IEEE
%% Document class options: replace "draftcls" by "final" for final document.
%% All other options may just as well be omitted because they are the default values.
%\documentclass[10pt,final,journal,letterpaper,twoside,twocolumn]{IEEEtran}


%%==========================================================================
%% Document automation
%%==========================================================================

\def\reptitle{Similarity measures for traffic scenarios}
\def\repauthor{Erwin de Gelder, etc.}


%%==========================================================================
%% Packages
%%==========================================================================

\usepackage[a4paper,left=3.5cm,right=3.5cm,top=3cm,bottom=3cm]{geometry} %% change page layout; remove for IEEE paper format
\usepackage[T1]{fontenc}                        %% output font encoding for international characters (e.g., accented)
\usepackage[cmex10]{amsmath}                    %% math typesetting; consider using the [cmex10] option
\usepackage{amssymb}                            %% special (symbol) fonts for math typesetting
\usepackage{amsthm}                             %% theorem styles
\usepackage{dsfont}                             %% double stroke roman fonts: the real numbers R: $\mathds{R}$
\usepackage{mathrsfs}                           %% formal script fonts: the Laplace transform L: $\mathscr{L}$
\usepackage[pdftex]{graphicx}                   %% graphics control; use dvips for TeXify; use pdftex for PDFTeXify
\usepackage{array}                              %% array functionality (array, tabular)
\usepackage{upgreek}                            %% upright Greek letters; add the prefix 'up', e.g. \upphi


\usepackage[utf8]{inputenc}    % utf8 support
\usepackage[style=ieee,doi=false,isbn=false,url=false,date=year,minbibnames=6,maxbibnames=15,backend=biber]{biblatex}
%\renewcommand*{\bibfont}{\footnotesize}
\setlength{\biblabelsep}{\labelsep}
\bibliography{../bib}

%\usepackage[round]{natbib}                     %% Author-year citations (remove package cite)
\usepackage{stfloats}                           %% improved handling of floats
\usepackage{multirow}                           %% cells spanning multiple rows in tables
%\usepackage{subfigure}                         %% subfigures and corresponding captions (for use with IEEEconf.cls)
%\usepackage{subfig}                             %% subfigures (IEEEtran.cls: set caption=false)
\usepackage{fancyhdr}                           %% page headers and footers
\usepackage[official,left]{eurosym}             %% the euro symbol; command: \euro
\usepackage{appendix}                           %% appendix layout
\usepackage{xspace}                             %% add space after macro depending on context
\usepackage{verbatim}                           %% provides the comment environment
\usepackage[dutch,USenglish]{babel}             %% language support
\usepackage{wrapfig}                            %% wrapping text around figures
\usepackage{longtable}                          %% tables spanning multiple pages
%\usepackage[utf8]{inputenc}
\usepackage{pgfplots}                           %% support for TikZ figures (Matlab)
\usepackage[breaklinks=true,hidelinks,          %% implement hyperlinks (dvips yields minor problems with breaklinks;
            bookmarksnumbered=true]{hyperref}   %% IEEEtran: set bookmarks=false)
%\usepackage[hyphenbreaks]{breakurl}            %% allow line breaks in URLs (don't use with PDFTeX)
\usepackage{lmodern} 
\usepackage{etoolbox}							%% Needed for apptocmd later
\usepackage[capitalize]{cleveref}
\crefname{figure}{Figure}{Figures}
\crefname{equation}{}{}
\usepackage{units}
\usepackage{subcaption}

%%==========================================================================
%% Fancy headers and footers
%%==========================================================================

\pagestyle{fancy}                                       %% set page style
\fancyhf{}                                              %% clear all header & footer fields
\fancyhead[L]{\includegraphics[width=15mm]{Streetwise}} %% define headers (LE: left field/even pages, etc.)
\fancyhead[R]{\small\emph{\reptitle}}                   %% similar
\fancyfoot[C]{\thepage}                                 %% define footer
\setlength{\headheight}{18pt}                           %% increase head height to accommodate an oversized picture in the header
\renewcommand*{\headrulewidth}{0.25pt}                  %% header line width
\renewcommand*{\footrulewidth}{0pt}                     %% footer line width
%% Redefine the default "plain" page style (automatically activated by \maketitle, \section, ...)
\fancypagestyle{plain}{\fancyhf{}
                       \fancyhead[C]{\includegraphics[width=30mm]{Streetwise}}
                       \renewcommand{\headrulewidth}{0pt}
                       \renewcommand{\footrulewidth}{0pt}
                       \setlength{\headheight}{36pt}}


%%==========================================================================
%% TikZ figures
%%==========================================================================

\newlength\figurewidth
\setlength\figurewidth{0.35\textwidth}              %% set figure height
\newlength\figureheight
\setlength\figureheight{0.3\textwidth}              %% set figure width
\pgfplotsset{every axis/.append style={
    scaled y ticks=false,
    scaled x ticks=false,
    y tick label style={/pgf/number format/fixed},
    x tick label style={/pgf/number format/fixed},
    legend style={font=\small}},
    compat=1.9}                                     %% PGFPlots package options
\usetikzlibrary{shapes.geometric, arrows, arrows.meta}
\usepgfplotslibrary{groupplots}
%\usetikzlibrary{external}                           %% Create pdf figures from TikZ. Use PDFTeXify ...
%\tikzexternalize[prefix=./tikz/]                    %% ... with --tex-option=--shell-escape switch.
%\tikzset{external/force remake}                    %% force pdf figure update


%%==========================================================================
%% User-defined commands
%%==========================================================================

\newcommand*{\mat}[1]{\mathbf{#1}}                              %% matrix/vector notation
\newcommand*{\matsym}[1]{\boldsymbol{#1}}                       %% matrix/vector notation for Greek letters
\newcommand*{\T}{^{\scriptscriptstyle\mathsf{T}}}               %% transpose operator
\newcommand*{\Hr}{^{\scriptscriptstyle\mathsf{H}}}              %% conjugate transpose operator
\newcommand*{\ud}{\mathrm{\,d}}                                 %% differential operator (upright d)
\newcommand*{\defeq}{\mathrel{\mathop:}=}                       %% definition sign :=
\newcommand*{\eqdef}{=\mathrel{\mathop:}}                       %% definition sign =:
\newcommand*{\ip}[2]{\left\langle#1\,{,}\,#2\right\rangle}      %% inner product
\newcommand*{\real}[1]{\mathrm{Re}(#1)}                         %% real part
\newcommand*{\imag}[1]{\mathrm{Im}(#1)}                         %% imaginary part
\newcommand*{\lsup}[1]{{}^{#1}\!}                               %% left superscript
\newcommand*{\hi}[1]{$^\text{#1}$}                              %% superscript in normal text
\newcommand*{\lo}[1]{$_\text{#1}$}                              %% subscript in normal text
\newcommand*{\w}[1]{\mathrm{#1}}                                %% multiple character super-/subscript in math mode
\newcommand*{\Ln}{{\ensuremath{\mathcal{L}}}\xspace}            %% L norm
\newcommand*{\Lone}{{\ensuremath{\mathcal{L}_1}}\xspace}        %% L_1 norm
\newcommand*{\Ltwo}{{\ensuremath{\mathcal{L}_2}}\xspace}        %% L_2 norm
\newcommand*{\Linf}{{\ensuremath{\mathcal{L}_\infty}}\xspace}   %% L_inf norm
\newcommand*{\Lp}{{\ensuremath{\mathcal{L}_p}}\xspace}          %% L_p norm
\newcommand*{\Htwo}{{\ensuremath{\mathcal{H}_2}}\xspace}        %% H_2 norm
\newcommand*{\Hinf}{{\ensuremath{\mathcal{H}_\infty}}\xspace}   %% H_inf norm
\newcommand*{\Hp}{{\ensuremath{\mathcal{H}_p}}\xspace}          %% H_p norm
\newcommand*{\capskip}{\vspace{-12pt}}                          %% caption skip for figures with subfloats
\newcommand*{\etal}{et al.}                                     %% may be required for Natbib bibliography styles
\renewcommand*{\qedsymbol}{$\blacksquare$}                      %% redefine the end-of-proof symbol
\renewcommand*{\labelitemi}{$\bullet$}                          %% first level item list bullet
\renewcommand*{\labelitemii}{$-$}                               %% second level item list bullet
%\renewcommand*{\theenumi}{\textit{\roman{enumi}}}               %% first level enumerator
\renewcommand*{\labelenumi}{\theenumi.}
\renewcommand*{\theenumii}{\textit{\alph{enumii}}}              %% second level enumerator
\renewcommand*{\labelenumii}{\theenumii.}
\DeclareMathOperator{\tr}{tr}                                   %% trace of a matrix
\DeclareMathOperator{\sgn}{sgn}                                 %% signum function
\DeclareMathOperator{\atan}{atan}                               %% arc tangent

\newcommand{\deuclid}{d_{\textup{Euclidian}}}
\newcommand{\imagi}{\textbf{j}}
\newcommand{\profile}[1]{\textbf{v}_{#1}}
\newcommand{\barprofile}[1]{\bar{\textbf{v}}_{#1}}


%%==========================================================================
%% User-defined environments
%%==========================================================================

\theoremstyle{plain}\newtheorem{definition}{Definition}[section]    %% definition
                    \newtheorem{theorem}{Theorem}[section]          %% theorem
                    \newtheorem{lemma}[theorem]{Lemma}              %% lemma
                    \newtheorem{corollary}[theorem]{Corollary}      %% corollary
                    \newtheorem{assumption}{Assumption}[section]    %% assumption
                    \newtheorem{condition}{Condition}[section]      %% condition
\theoremstyle{definition}\newtheorem{example}{Example}[section]     %% examples
\theoremstyle{remark}\newtheorem{remarkenv}{Remark}[section]        %% remarks
\newenvironment{remark}{\begin{remarkenv}}%
                       {\hfill$\blacklozenge$\end{remarkenv}}       %% end remark with a lozenge
\newenvironment{reviewer}{\itshape}{\upshape}                       %% environment for reviewer's comments


%%==========================================================================
%% Miscellaneous
%%==========================================================================

\graphicspath{{./../}{./figures/}{./more_figures/}} %% (graphicx) directory path for figures
%\setlength{\parindent}{0pt}                        %% no paragraph indentation
%\setlength{\parskip}{2ex}                          %% create empty line between paragraphs
%\interdisplaylinepenalty=2500                      %% (amsmath) allow for page breaks within multiline equations
%\numberwithin{equation}{section}                   %% (amsmath) include section number in equation numbering
%\numberwithin{figure}{section}                     %% (amsmath) include section number in figure numbering
%\numberwithin{table}{section}                      %% (amsmath) include section number in table numbering
\addtolength{\arraycolsep}{-0.5mm}                  %% squeeze matrix columns a little
\fnbelowfloat                                       %% (stfloats) put footnote below a float at the page bottom
\urlstyle{same}                                     %% (hyperref) use current font for URLs
\hypersetup{pdftitle={\reptitle},
            pdfauthor={\repauthor}}                 %% (hyperref) pdf properties title and author
%\raggedbottom                                      %% don't add inter-paragraph spacing to achieve \textheight
%\setlength\subfigcapskip{-3pt}                     %% (subfigure) distance between subfloat and subcaption
%\setlength\subfigbottomskip{-3pt}                  %% (subfigure) distance between subcaption and caption
%\renewcommand*{\subcapsize}{\small}                %% (subfigure) subcaption font size
\renewcommand*{\thesubfigure}{(\alph{subfigure})}   %% (subfig) implement 1(a) instead of 1a ...as subfigure reference
\captionsetup[subfloat]{labelformat=simple}        %% ... as subfigure reference
%\captionsetup[subfloat]{
%    farskip=0pt,
%    nearskip=8pt,
%    captionskip=1.5pt,
%    labelfont={small,bf},
%    textfont=small}                                %% (subfig) subfloat caption format
%\captionsetup[figure]{
%    labelfont={small,bf},
%    textfont=small}                                %% (subfig/caption) figure caption format
%\captionsetup[table]{
%    aboveskip=2pt,
%    labelfont={small,bf},
%    textfont={small,sc}}                           %% (subfig/caption) table caption format
\apptocmd{\thebibliography}{\raggedright}{}{}		%% Suppress badness warnings in bibliography


% Table stuff
\usepackage{booktabs}
\usepackage{tabularx}
\setlength{\heavyrulewidth}{0.1em}
\newcommand{\otoprule}{\midrule[\heavyrulewidth]}

%%==========================================================================
%% Begin document
%%==========================================================================

\begin{document}

\selectlanguage{USenglish}

\title{\textbf{\reptitle}}
\author{\repauthor}
\date{\today}
\maketitle

\tableofcontents

\newpage

\section{Introduction}
\label{sec:introduction}
% Scenario-based approach for assessment of AVs -> this scenarios are important
An important aspect in the development of automated vehicles (AVs) is the assessment of quality and performance aspects of the AVs, such as safety, comfort, and efficiency \cite{bengler2014threedecades, stellet2015taxonomy, putz2017pegasus, roesener2016scenariobased, wachenfeld2016release, kompass2015sicherheitsveranderung}. 
For legal and public acceptance, it is important that there is a clear definition of system performance and that there are quantitative measures for the system quality. 
The more traditional methods \cite{response2006code, ISO26262}, used for evaluation of driver assistance systems, are no longer valid for the assessment of quality and performance aspects of an AV \cite{wachenfeld2016release}. 
Therefore, a scenario-based approach is proposed \cite{roesener2016scenariobased, putz2017pegasus, kompass2015sicherheitsveranderung}. 
For the scenario-based assessment, proper specification of scenarios is crucial since they are directly reflected in test cases used for scenario-based assessment \cite{stellet2015taxonomy}.

% Purpose of document
To structure the variety of scenarios, the notion of scenarios classes is introduced \cite{degelder2018ontology}. Whereas a scenario is a quantitative description, a scenario class is a qualitative description \cite{degelder2018ontology}. This document describes different measures to quantify the (dis)similarity between two different scenarios that belong to the same scenario class. An appropriate (dis)similarity measure for scenarios would be useful for at least two reasons.

% Use of a measure for completeness
The first use of a (dis)similarity measure is for quantifying the so-called \emph{completeness}. In order to draw conclusions on how an automated vehicle would perform in real-world traffic, it is necessary to know how representative the scenario catalog, which is used for the scenario-based assessment of the automated vehicle, is. Therefore it is important to quantify how complete the scenario catalog is \cite{geyer2014, alvarez2017prospective, stellet2015taxonomy}. A (dis)similarity measure of scenarios is useful for quantifying the completeness. For example, if newly collected scenarios are not similar to earlier collected scenarios, it is an indicated that the set of collected scenarios is not \emph{complete}. To achieve a complete set of scenarios, this process could be repeated until the newly collected scenarios are - to a certain extent - similar to the early collected scenarios. This process of achieving a complete set of scenarios is visualized in \cref{fig:completeness collection}.

\begin{figure}[b]
	\centering
	\tikzstyle{startstop}=[rectangle, rounded corners, minimum width=4em, minimum height=4em, text centered, draw=black, fill=red!30, text width=4em]
	\tikzstyle{process} = [rectangle, minimum width=9em, text width=8.5em, minimum height=4em, text centered, draw=black, fill=orange!30, node distance=10em]
	\tikzstyle{decision} = [diamond, minimum width=4em, text width=4em, minimum height=4em, text centered, draw=black, fill=green!30, node distance=10em]
	\tikzstyle{arrow} = [-{Triangle[width=2mm, length=3mm, round]}, line width=0.6mm] % rounded corners=0.4em
	\begin{tikzpicture}
		% Draw the blocks
		\node[startstop](start){Start};
		\node[process, right of=start, node distance=7em, text width=5em, minimum width=5.5em](collect){Collect new scenarios};
		\node[process, right of=collect, node distance=9em](compare){Compare new scenarios with scenario database};
		\node[decision, right of=compare](similar){New scenarios similar?};
		\node[startstop, right of=similar, node distance=9em](stop){Scenario database complete};
		\node[process, below of=compare, node distance=7em](add){Add collected scenarios to scenario database};
		
		% Arrows
		\draw[arrow] (start) -- (collect);
		\draw[arrow] (collect) -- (compare);
		\draw[arrow] (compare) -- (similar);
		\draw[arrow] (similar) -- node[anchor=south]{Yes} (stop);
		\draw[arrow] (similar) |- node[anchor=north]{No} (add);
		\draw[arrow] (add) -| (collect);
	\end{tikzpicture}
	\caption{Schematic process of collecting scenarios of a specific \emph{scenario class} until a certain completeness is reached.}
	\label{fig:completeness collection}
\end{figure}

% Use of a measure for scenario generation
The second use of a (dis)similarity measure is for the generation of test scenarios. Using collected real-world scenarios, one might want to generate an additional set of scenarios to be used for the scenario-based assessment in order to emphasize safety-critical or performance-critical scenarios (cf. \cite{zofka2015datadrivetrafficscenarios, deGelder2017assessment}). Although the generated scenarios are different from the collected scenarios, it is still desirable that the generated scenarios are reflecting what might be encountered in real-life. For this purpose, a (dis)similarity measure can be used. For example, if the generated scenarios is similar to one or several encountered scenarios, it is more conceivable that this scenario is a realistic scenario compared to a scenario that does not show high similarities with any of the collected scenarios. 


\section{Dissimilarity measures}
\label{sec:dissimilarity}

In this section, several dissimilarity measures are presented. They can be divided into four different categories \cite{serra2014empirical}: lock-step measures, feature-based measures, model-based measures, and elastic measures. Before discussing these different categories, some necessary notation details are described.

\subsection{Notation and example}
\label{sec:dissimilarity notation}

It is assumed that the scenario is described by means of a time series. Therefore, the dissimilarity measures compare the time series $\textbf{x}$ and $\textbf{y}$. The i-th element of $\textbf{x}$ and $\textbf{y}$ are denoted by $x_i$ and $y_i$, respectively. For the sake of simplicity, $x_i$ and $y_i$ are scalars. The length of the time series $\textbf{x}$ and $\textbf{y}$ is $N$ and $M$, respectively. Note that the time series do not necessarily have the same length, i.e., $N$ might be different than $M$. It is assumed, however, that the sample time is constant and similar for both time series.

To illustrate the dissimilarity measures, four different velocity profiles are considered, see \cref{fig:example time series}. All four velocity profiles describe the velocity of a vehicle while braking. The average deceleration is $\unit[3.19]{m/s^2}$, $\unit[1.72]{m/s^2}$, $\unit[1.63]{m/s^2}$, and $\unit[1.64]{m/s^2}$ respectively. We will refer to the the profiles with $\profile{1}$, $\profile{2}$, $\profile{3}$, and $\profile{4}$, respectively. 

\begin{figure}
	\centering
	\setlength\figureheight{0.4\linewidth}
	\setlength\figurewidth{0.7\linewidth}
	% This file was created by matplotlib2tikz v0.6.14.
\begin{tikzpicture}

\definecolor{color0}{rgb}{0.12156862745098,0.466666666666667,0.705882352941177}

\begin{axis}[
xlabel={Time [s]},
ylabel={Velocity [km/h]},
xmin=-0.293991395216091, xmax=6.17381929953791,
ymin=7.506, ymax=35.622,
width=\figurewidth,
height=\figureheight,
tick align=outside,
tick pos=left,
xmajorgrids,
x grid style={white!69.019607843137251!black},
ymajorgrids,
y grid style={white!69.019607843137251!black}
]
\addplot [semithick, color0, forget plot]
table {%
0 33.984
0.0399989420671147 34.344
0.0799978841342437 34.02
0.119996826201358 34.056
0.159996163019557 34.02
0.199995105086685 34.128
0.2399940471538 34.056
0.279992989220929 34.02
0.319991931288044 33.84
0.359990873355173 33.876
0.399989815422288 33.768
0.439988757489417 33.588
0.479987699556531 33.408
0.51998664162366 33.264
0.559985583690775 33.012
0.599997157792259 32.832
0.679995041926503 32.544
0.719993983993632 32.292
0.759992926060747 32.22
0.799991868127876 31.86
0.83999081019499 31.644
0.879989752262119 31.536
0.919988694329234 31.248
0.959987636396363 30.852
0.999986578463478 30.672
1.03998552053061 30.492
1.07998446259772 30.132
1.11998340466485 29.952
1.15998234673197 29.592
1.19998128879909 29.304
1.23998023086621 29.088
1.27997917293334 28.836
1.31997811500045 28.584
1.35997705706758 28.296
1.3999759991347 27.864
1.43997494120183 27.72
1.47997427802002 27.324
1.51997322008714 27.036
1.55997216215427 26.784
1.59997110422138 26.46
1.63997004628851 26.136
1.67996898835563 25.92
1.71995529838838 25.632
1.75995424045551 25.2
1.79995318252263 25.092
1.83995172983867 24.732
1.87995106665687 24.372
1.919950008724 24.084
1.95994895079112 23.832
1.99994789285824 23.472
2.03994683492536 23.184
2.07994577699249 22.896
2.1199447190596 22.536
2.15994366112673 22.356
2.19994260319385 22.068
2.23994154526098 21.6
2.31993942939522 21.204
2.35993837146233 20.808
2.39993731352946 20.592
2.43993625559658 20.196
2.47993519766371 19.944
2.51993413973082 19.8
2.55993308179795 19.512
2.59993202386507 19.224
2.63993096593219 19.044
2.67992990799931 18.756
2.71992885006644 18.468
2.75992779213355 18.324
2.79992673420068 18.036
2.8399256762678 17.784
2.87992461833493 17.712
2.91992356040204 17.496
2.95992250246917 17.28
3.03992038660341 17.1
3.07991932867053 16.992
3.11991827073766 17.028
3.15991721280477 16.848
3.23990364915787 16.776
3.279902591225 16.668
3.31990153329211 16.596
3.35990047535924 16.524
3.39989941742635 16.38
3.4798973015606 16.344
3.51989624362773 16.2
3.59989412776197 16.092
3.63989306982909 16.02
3.67989201189621 15.984
3.71989095396333 15.948
3.75988989603046 15.804
3.8398877801647 15.624
3.87988672223183 15.48
3.95988460636607 15.372
3.99988354843319 15.156
4.03988249050032 15.192
4.07988143256743 14.976
4.11988037463456 14.832
4.15987931670168 14.652
4.19987825876881 14.544
4.23987720083592 14.436
4.27987614290305 14.4
4.31987508497016 14.22
4.35987402703729 14.04
4.39987296910441 14.004
4.43987191117154 13.788
4.47987085323865 13.644
4.51986979530578 13.5
4.55986873737289 13.32
4.59986767944002 13.176
4.63986662150714 13.068
4.67986556357427 12.996
4.71986450564138 12.744
4.75986344770851 12.636
4.79986238977563 12.42
4.83986133184276 12.276
4.87986027390987 12.204
4.919859215977 12.096
4.95985815804411 11.988
4.99985710011124 11.808
5.03985012091225 11.592
5.07984906297938 11.52
5.11984800504649 11.412
5.15984694711362 11.196
5.19984588918074 11.124
5.23984483124786 10.944
5.27984377331498 10.764
5.31984271538211 10.692
5.35984165744922 10.548
5.39984059951635 10.332
5.43983954158347 10.26
5.4798384836506 10.044
5.51983742571771 9.864
5.55983636778484 9.792
5.59983530985195 9.612
5.63983425191908 9.468
5.6798331939862 9.36
5.71983213605333 9.216
5.75983107812044 9.072
5.79983002018757 9
5.83982896225469 8.82
5.87982790432181 8.784
};
\end{axis}

\end{tikzpicture}
	\caption{Four examples of time series which will be used to illustrate the dissimilarity measures. The velocity is sampled with $\unit[10]{Hz}$.}
	\label{fig:example time series}
\end{figure}

\subsection{Lock-step measures}
\label{sec:lock step}

Lock-step measures estimate the dissimilarity between two time series by comparing each time instant. This, however, does only work if the two time series have equal length. In case that the time series $\textbf{x}$ and $\textbf{y}$ do not have the same length, one of the time series can be re-sampled. This works well for a number of data sources \cite{keogh2003survey}. Let the time series $\bar{\textbf{y}}$ denote the re-sampled time series $\textbf{y}$, such that $\bar{\textbf{y}}$ contains $N$ samples, e.g., using linear interpolation. As a consequence, the elements $\bar{y}[i]$ of $\bar{\textbf{y}}$ are not sampled with the same frequency as the time series $\textbf{x}$ and $\textbf{y}$. 

A straightforward way to estimate the dissimilarity between two time series is to use an $L_p$ norm. The $L_p$ norm is as follows:
\begin{equation} \label{eq:lp norm}
	d_{L_p}(\textbf{x}, \bar{\textbf{y}}) = \left( \sum_{i=0}^{N-1} |x[i] - \bar{y}[i]|^p \right)^{\frac{1}{p}},
\end{equation}
where $p \geq 0$. Not that with $p=2$, we obtain the Euclidean distance:
\begin{equation} \label{eq:euclidian}
	\deuclid(\textbf{x}, \bar{\textbf{y}}) = \sqrt{\sum_{i=0}^{N-1} |x[i] - \bar{y}[i]|^p}.
\end{equation}
The Euclidean distance is often applied due to its computational simplicity \cite{xing2010survey}. Even if there might be better dissimilarity measures for certain applications, the Euclidean distance serves as a good benchmark \cite{keogh2003survey}.

For our purposes, the Euclidean distances is not suitable. Because, in general, we have $N \ne M$, we need to re-sample one time series. Because of this, a strong braking action might look `similar' to very smooth braking, see \cref{fig:resampled profiles}. In \cref{fig:resampled profiles}, $\textbf{v}_1$ and $\textbf{v}_3$ are re-sample, such that they have the same amount of samples as $\profile{2}$, resulting in $\barprofile{1}$ and $\barprofile{3}$, respectively. This results in $\deuclid(\profile{2}, \barprofile{1})=\unit[13.3]{m/s}$ and $\deuclid(\profile{2}, \barprofile{3})=\unit[14.7]{m/s}$. So even though $\profile{3}$ looks more similar to $\profile{2}$ than $\profile{1}$, the Euclidean distance \cref{eq:euclidian} gives a higher score for $\deuclid(\profile{2}, \barprofile{3})$ compared to $\deuclid(\profile{2}, \barprofile{1})$.

\begin{figure}
	\centering
	\setlength\figureheight{0.4\linewidth}
	\setlength\figurewidth{0.46\linewidth}
	\begin{minipage}{0.46\linewidth}
		\centering
		% This file was created by matplotlib2tikz v0.6.14.
\begin{tikzpicture}

\definecolor{color1}{rgb}{1,0.498039215686275,0.0549019607843137}
\definecolor{color0}{rgb}{0.12156862745098,0.466666666666667,0.705882352941177}

\begin{axis}[
xlabel={Sample},
ylabel={Velocity [km/h]},
xmin=-3.8, xmax=79.8,
ymin=-2.4336, ymax=51.1056,
width=\figurewidth,
height=\figureheight,
tick align=outside,
tick pos=left,
xmajorgrids,
x grid style={lightgray!92.026143790849673!black},
ymajorgrids,
y grid style={lightgray!92.026143790849673!black}
]
\addplot [semithick, color0, mark=*, mark size=1, mark options={solid}, only marks, forget plot]
table {%
0 47.088
1 46.9079863570422
2 46.7999676690327
3 46.5479802461338
4 46.3679702254707
5 46.1159598531504
6 45.9358601596821
7 45.4498142231174
8 45.179940454173
9 44.729704544694
10 44.3158188044651
11 43.8836856435103
12 43.3796156522959
13 42.767789955294
14 42.0838333001078
15 41.5075281810371
16 40.8955037374065
17 40.2294110079255
18 39.5993835088374
19 38.8795685414562
20 37.9797001391539
21 37.3852946164394
22 36.68367427438
23 35.9631533633361
24 35.3150331264124
25 34.5051798318338
26 33.6233394534536
27 32.9032192325912
28 32.111495850071
29 31.4444417183995
30 30.4553562229219
31 29.9144526054014
32 29.2667486138953
33 28.3131804140041
34 27.5397595428659
35 26.8006419351036
36 26.2428597471555
37 25.5592204405501
38 24.8025352295064
39 24.0298637835918
40 23.3986030116011
41 22.6793546660349
42 21.9596708636399
43 21.4004881897229
44 20.7706317250088
45 20.2666044390949
46 19.6907574960284
47 19.1152763547719
48 18.5036320664783
49 17.9985009799961
50 17.5308574022717
51 17.0268369378363
52 16.4870142531641
53 15.8930926011825
54 15.299260445287
55 14.7406843782999
56 14.1466629901203
57 13.5526391149456
58 12.958617726766
59 12.453991216557
60 11.8423685296134
61 11.1595853109241
62 10.4393688000234
63 9.80807570341702
64 9.21361447195764
65 8.53068047016063
66 7.84599002842424
67 7.10886806083318
68 6.44308081477942
69 5.81190431375665
70 5.14634863363259
71 4.56960679779292
72 3.92181861892924
73 3.2390168543433
74 2.51875579065799
75 1.94143128682471
76 0
};
\addplot [semithick, color1, mark=*, mark size=1, mark options={solid}, only marks, forget plot]
table {%
0 48.672
1 48.6511532049701
2 48.5223070494538
3 48.1054625993072
4 47.7454535228184
5 47.432807257711
6 46.9496484429941
7 46.4266970766413
8 45.884736970841
9 45.3352157194084
10 44.7724624631194
11 44.1433816638767
12 43.4763994115175
13 42.8283598893067
14 42.1935904806297
15 41.6100123542878
16 40.9506358542192
17 40.2003013056784
18 39.5182638690388
19 38.8618378493623
20 38.1010938482656
21 37.3772736328982
22 36.7519168461425
23 36.058390864896
24 35.3080572215737
25 34.6173688270732
26 33.9418322161786
27 33.2852684439624
28 32.5605748812239
29 31.7166541850711
30 30.8858475066327
31 30.0625345526024
32 29.3338009525241
33 28.6107804682413
34 27.8709082957595
35 27.1405119392004
36 26.4214865617486
37 25.6872346559076
38 24.9472728019208
39 24.0929591732244
40 23.2953374654076
41 22.6488942132699
42 21.8093904457685
43 20.8090029154641
44 19.9962709776019
45 19.2214236767378
46 18.4084103182054
47 17.6030181206217
48 16.8109629546984
49 15.9989914894015
50 15.1756392816052
51 14.3902676830082
52 13.5974031316081
53 12.7537290184546
54 11.9288019127081
55 11.12637121585
56 10.4525783200168
57 9.8270245995681
58 9.22262864130256
59 8.63526267358882
60 8.0933100140132
61 7.54001157441209
62 6.97725831812309
63 6.37213018030013
64 5.75944392605043
65 5.16534291094228
66 4.58636348414682
67 4.03384683689853
68 3.52126880464978
69 3.03151283811506
70 2.57031470907598
71 2.12298857589045
72 1.72523539717292
73 1.1113635605376
74 0.23814933440091
75 0
76 0
};
\end{axis}

\end{tikzpicture}
		\subcaption{$\deuclid(\textbf{v}_2, \bar{\textbf{v}}_1)=\unit[13.3]{m/s}$, $\textbf{v}_2$ in blue and $\bar{\textbf{v}}_1$ in orange.}
	\end{minipage}
	\hspace{0.04\linewidth}
	\begin{minipage}{0.46\linewidth}
		% This file was created by matplotlib2tikz v0.6.14.
\begin{tikzpicture}

\definecolor{color1}{rgb}{1,0.498039215686275,0.0549019607843137}
\definecolor{color0}{rgb}{0.12156862745098,0.466666666666667,0.705882352941177}

\begin{axis}[
xlabel={Sample},
ylabel={Velocity [km/h]},
xmin=-3.7, xmax=77.7,
ymin=0.0601814252858692, ymax=50.9868485035578,
width=\figurewidth,
height=\figureheight,
tick align=outside,
tick pos=left,
xmajorgrids,
x grid style={white!69.019607843137251!black},
ymajorgrids,
y grid style={white!69.019607843137251!black}
]
\addplot [semithick, color0, mark=*, mark size=1, mark options={solid}, only marks, forget plot]
table {%
0 47.088
1 46.9079863570422
2 46.7999676690327
3 46.5479802461338
4 46.3679702254707
5 46.1159598531504
6 45.9358601596821
7 45.4498142231174
8 45.179940454173
9 44.729704544694
10 44.3158188044651
11 43.8836856435103
12 43.3796156522959
13 42.767789955294
14 42.0838333001078
15 41.5075281810371
16 40.8955037374065
17 40.2294110079255
18 39.5993835088374
19 38.8795685414562
20 37.9797001391539
21 37.3852946164394
22 36.68367427438
23 35.9631533633361
24 35.3150331264124
25 34.5051798318338
26 33.6233394534536
27 32.9032192325912
28 32.111495850071
29 31.4444417183995
30 30.4553562229219
31 29.9144526054014
32 29.2667486138953
33 28.3131804140041
34 27.5397595428659
35 26.8006419351036
36 26.2428597471555
37 25.5592204405501
38 24.8025352295064
39 24.0298637835918
40 23.3986030116011
41 22.6793546660349
42 21.9596708636399
43 21.4004881897229
44 20.7706317250088
45 20.2666044390949
46 19.6907574960284
47 19.1152763547719
48 18.5036320664783
49 17.9985009799961
50 17.5308574022717
51 17.0268369378363
52 16.4870142531641
53 15.8930926011825
54 15.299260445287
55 14.7406843782999
56 14.1466629901203
57 13.5526391149456
58 12.958617726766
59 12.453991216557
60 11.8423685296134
61 11.1595853109241
62 10.4393688000234
63 9.80807570341702
64 9.21361447195764
65 8.53068047016063
66 7.84599002842424
67 7.10886806083318
68 6.44308081477942
69 5.81190431375665
70 5.14634863363259
71 4.56960679779292
72 3.92181861892924
73 3.2390168543433
74 2.51875579065799
};
\addplot [semithick, color1, mark=*, mark size=1, mark options={solid}, only marks, forget plot]
table {%
0 48.672
1 48.592454630678
2 48.3898238011854
3 48.1090899735643
4 47.8916293771907
5 47.5525254382995
6 47.2017787388955
7 46.7609227792633
8 46.2200341165703
9 45.6614714914681
10 45.183685772402
11 44.7380017794322
12 44.2592791562179
13 43.8244831141074
14 43.3545809807361
15 42.8007607667604
16 42.3083539476533
17 41.7392054328171
18 41.0385345936437
19 40.5026559154203
20 39.821503156755
21 39.1327865527472
22 38.4403704115835
23 37.7350266409662
24 37.043224156655
25 36.2800091920934
26 35.4146260181337
27 34.6493599816715
28 33.7692810508542
29 33.0988640948658
30 32.4457548748638
31 31.9006367617273
32 30.8999178027883
33 30.1234484003732
34 29.4538699781143
35 28.7593308876634
36 28.2017371143271
37 27.5027206047803
38 26.8728414941303
39 26.3065229665031
40 25.7317558816128
41 25.1661662550748
42 24.5372038663523
43 23.9656126536124
44 23.3591586262176
45 22.6904438245024
46 22.1340478255719
47 21.627202202576
48 21.0353364672983
49 20.4503075947583
50 19.7923859917776
51 19.1616198509756
52 18.5629825152108
53 17.9122437313713
54 17.2637963699814
55 16.5750518228492
56 15.8458432785579
57 15.2300895095358
58 14.535958056179
59 13.8899363260535
60 13.2448055159911
61 12.6019731940591
62 11.8915397692451
63 11.1445351298071
64 10.3857612631612
65 9.62460944722811
66 8.75722915569906
67 7.82468199937867
68 6.90767196831181
69 5.98131586183445
70 5.06367145219246
71 4.16611312755441
72 3.35315696512353
73 2.46436354889437
74 2.37502992884368
};
\end{axis}

\end{tikzpicture}
		\subcaption{$\deuclid(\textbf{v}_2, \bar{\textbf{v}}_3)=\unit[14.7]{m/s}$, $\textbf{v}_2$ in blue and $\bar{\textbf{v}}_3$ in orange.}
	\end{minipage}
	\caption{Comparison of the different profiles, where one of the two profiles is re-sampled.}
	\label{fig:resampled profiles}
\end{figure}

\subsection{Feature-based measures}
\label{sec:feature based}

The Euclidean distance of \cref{eq:euclidian} is not of direct use for our application. An simple extension, however, is useful. This extension is to compute the Euclidean distance not directly using the (re-sampled) time series, but using features extracted from it \cite{serra2014empirical}. Let $\hat{x}_k$ and $\hat{y}_k$ with $k \in \{1, \ldots, \theta\}$ denote features extracted from the time series $\textbf{x}$ and $\textbf{y}$, respectively. The dissimilarity can be formulated as follows:

\begin{equation}\label{eq:eucl features}
	d_{\textup{F}}(\textbf{x}, \textbf{y}) = \sqrt{ \sum_{k=1}^{\theta} w_k \left( \hat{x}_k - \hat{y}_k \right)^2 }.
\end{equation}
Here, $w_k$ denote weights that can be chosen to emphasize the $k$-th parameter less or more. 

\subsubsection{Fourier coefficients}
\label{sec:fourier coefficient}
One possibility is to use the Fourier coefficients as parameters \cite{agrawal1993similarity, faloutsos1994fast}. The Fourier coefficients are obtained with the Discrete Fourier Transform \cite{oppenheim1999discrete}. The $n$-th Fourier coefficient, i.e.,
\begin{equation}
	X[n] = \sum_{i=0}^{N-1} x[i] e^{-\imagi\frac{2\pi i n}{N}},
\end{equation}
correspond to a frequency of $f_n = n/Nt_s$, where $N$ is the total number of samples and $t_s$ is the sample time. When only the Fourier coefficients corresponding to the lower frequencies are used as parameters, this automatically has the effect of removing rapidly-fluctuating components of the time series \cite{serra2014empirical}.

Note that the Fourier coefficients of different time series correspond to different frequencies, due to the dependency of the number of samples. In order to compare the different time series of \cref{fig:example time series}, linear interpolation is used to retrieve the Fourier coefficients at $\unit[1]{Hz}$, $\unit[2]{Hz}$, and $\unit[3]{Hz}$. Furthermore, the resulting coefficients are divided over the number of samples, see \cref{tab:fourier coefficients}.

\begin{table}
	\centering
	\caption{Interpolated Fourier coefficients for the different time series (see \cref{fig:example time series}) divided over the number of samples.}
	\label{tab:fourier coefficients}
	\begin{tabular}{crrrr}
    \toprule
    Time    & \multicolumn{3}{c}{Fourier coefficients} \\
    series  & $\unit[0.0]{Hz}$ & $\unit[0.5]{Hz}$ & $\unit[1.0]{Hz}$ & $\unit[1.5]{Hz}$ \\\otoprule
    $\profile{1}$ & $25.15+0.00\imagi$ & $0.60-3.76\imagi$ & $0.52-1.78\imagi$ & $0.53-1.11\imagi$ \\
    $\profile{2}$ & $25.52+0.00\imagi$ & $0.21-1.94\imagi$ & $0.27-0.94\imagi$ & $0.28-0.61\imagi$ \\
    $\profile{3}$ & $27.37+0.00\imagi$ & $0.20-1.90\imagi$ & $0.25-0.89\imagi$ & $0.28-0.58\imagi$ \\
    $\profile{4}$ & $23.90+0.00\imagi$ & $0.19-2.38\imagi$ & $0.27-0.99\imagi$ & $0.30-0.59\imagi$ \\
    \bottomrule
\end{tabular}

\end{table}

\subsubsection{Auto-regression coefficients}
\label{sec:auto regression coefficients}
Instead of using the Fourier coefficients as parameters, one could use the coefficients from an auto-regressive (AR) model. An AR model approximates the time series using the model 
\begin{equation}
	\label{eq:auto-regression}
	\hat{x}[i] = a_0 + \sum_{j=1}^{\eta-1} a_j \hat{x}[i-j],
\end{equation}
where $\eta$ denotes the number of coefficients. 

\begin{table}
	\centering
	\caption{Coefficients of auto-regression model \cref{eq:auto-regression} with $\eta=3$ for the three time series of \cref{fig:example time series}.}
	\label{tab:ar coefficients}
	\begin{tabular}{crrr}
    \toprule
    Time    & \multicolumn{3}{c}{Auto-regression coefficients} \\
    series & $a_0\,[\unit{m/s}]$ & $a_1$ & $a_2$ \\ \otoprule
    $\profile{1}$ & -0.13 & 1.57 & -0.57 \\
    $\profile{2}$ & -0.06 & 1.66 & -0.66 \\
    $\profile{3}$ & -0.12 & 1.43 & -0.43 \\
    $\profile{4}$ & -0.03 & 1.77 & -0.77 \\
    \bottomrule
\end{tabular}

\end{table}

\subsubsection{Custom scenario features}
\label{sec:custom features}

\begin{table}
	\centering
	\caption{Feature-based dissimilarity measures, according to \cref{eq:eucl features}, with Fourier coefficients (see \cref{sec:fourier coefficient}), auto-regression coefficients (see \cref{sec:auto regression coefficients}), and custom scenario features (see \cref{sec:custom features}).}
	\label{tab:feature based dissimilarity}
	\begin{tabular}{cccccc}
    \toprule
    \multicolumn{2}{c}{Time series} & \multicolumn{2}{c}{Features based} & \multicolumn{2}{c}{Model based} \\
    \textbf{x} & \textbf{y} & Fourier coefficients & Custom features & Auto regression & Splines \\ \otoprule
    $\profile{1}$ & $\profile{2}$ & $5.22$ & $1.01$ & $0.53$ & $0.23$ \\
    $\profile{1}$ & $\profile{3}$ & $5.49$ & $1.22$ & $1.13$ & $0.25$ \\
    $\profile{1}$ & $\profile{4}$ & $5.03$ & $2.41$ & $2.67$ & $2.86$ \\
    $\profile{2}$ & $\profile{3}$ & $0.30$ & $0.24$ & $1.52$ & $0.07$ \\
    $\profile{2}$ & $\profile{4}$ & $0.46$ & $2.75$ & $2.16$ & $2.46$ \\
    $\profile{3}$ & $\profile{4}$ & $0.63$ & $2.95$ & $3.64$ & $2.60$ \\
    \bottomrule
\end{tabular}

\end{table}

\printbibliography

\end{document} 